% 075_RSA_En_ch.tex -- chapter content Doc. 075 EN
% RSA and Period Finding: Foundations, Resonance Formulation, Limits
% Unified version August 2026 (R75)

\section*{Summary}

This document treats the RSA procedure in two parts. Part~I presents
the number-theoretic foundations and security analysis. Part~II
examines the attack surface via period finding as a resonance problem
and determines its carriers and limits: the resonance description is
correct -- period finding is a spectral problem -- but its generators
are the \textbf{primes}, not an adjustable parameter.

\begin{keyresult}[Two separate parameters]
This document uses two quantities that must not be confused:
\begin{itemize}
  \item $\xi = 4/30000$ -- the FFGFT geometry parameter. Fixed, not
    adjustable, with no operational connection to factorisation.
  \item $\sigma$ -- a numerical grid parameter of the search algorithms
    (values such as $1/42$, $1/100$, $1/1000$). A sampling quantity,
    not a physical one.
\end{itemize}
\end{keyresult}

%-------------------------------------------------------------------
\part*{Part I: The RSA Procedure}
%-------------------------------------------------------------------

\section{Introduction to RSA Cryptography}

The RSA algorithm (Rivest, Shamir, Adleman 1977) is one of the first
practical public-key cryptosystems and is widely used for secure data
transmission.

\subsection{Mathematical Foundation}

RSA is based on the computational difficulty of factoring large integers
and the properties of modular arithmetic.

\subsection{Key Generation}

\begin{enumerate}
  \item Choose two distinct prime numbers $p$ and $q$.
  \item Compute $n = p \cdot q$.
  \item Compute Euler's totient: $\varphi(n) = (p-1)(q-1)$.
  \item Choose $e$ with $1 < e < \varphi(n)$ and $\gcd(e,\varphi(n))=1$.
  \item Compute $d$ with $d \cdot e \equiv 1 \pmod{\varphi(n)}$.
\end{enumerate}

The public key is $(n,e)$, the private key $(n,d)$.

\subsection{Encryption and Decryption}

For a message $M$ with $0 \le M < n$:
\begin{align}
  C &\equiv M^e \pmod n \qquad \text{(encryption)}, \\
  M &\equiv C^d \pmod n \qquad \text{(decryption)}.
\end{align}

\section{Mathematical Proofs}

\subsection{Correctness Proof}

\begin{theorem}[RSA Correctness]
For any message $M$ with $0 \le M < n$ and $\gcd(M,n)=1$:
$(M^e)^d \equiv M \pmod n$.
\end{theorem}
\begin{proof}
From $ed \equiv 1 \pmod{\varphi(n)}$ follows $ed = 1+k\varphi(n)$.
By Euler's theorem $M^{\varphi(n)} \equiv 1 \pmod n$, so
\[
  C^d \equiv M^{ed} \equiv M^{1+k\varphi(n)} \equiv
  M\cdot(M^{\varphi(n)})^k \equiv M \pmod n.
\]
For $\gcd(M,n)\neq 1$ the Chinese Remainder Theorem gives the same.
\end{proof}

\section{Implementation}

\subsection{Modular Exponentiation}

Efficient modular exponentiation by square-and-multiply:

\begin{tcolorbox}[colback=blue!5!white,colframe=blue!75!black,
  title=Algorithm: Modular Exponentiation]
\textbf{Function} ModExp($base$, $exponent$, $modulus$):
\begin{enumerate}
  \item $result \gets 1$; $base \gets base \bmod modulus$
  \item \textbf{while} $exponent > 0$:
  \begin{enumerate}
    \item \textbf{if} $exponent \bmod 2 = 1$:
      $result \gets (result \times base) \bmod modulus$
    \item $exponent \gets \lfloor exponent/2 \rfloor$;
      $base \gets (base \times base) \bmod modulus$
  \end{enumerate}
  \item \textbf{return} $result$
\end{enumerate}
\end{tcolorbox}

\subsection{Prime Generation}

\begin{itemize}
  \item Probabilistic primality tests (Miller-Rabin).
  \item $p$ and $q$ of similar bit length.
  \item Avoid special forms that are easier to factor.
\end{itemize}

\section{Security Analysis}

\subsection{Common Attacks}

\begin{center}
\begin{tabular}{ll}
\toprule
Attack & Description \\
\midrule
Factorisation & Factor $n$ into $p$ and $q$ \\
Small $e$ & Certain messages recoverable \\
Timing attacks & Deduce secret from computation time \\
Side-channel & Power, electromagnetic leaks \\
\bottomrule
\end{tabular}
\end{center}

\subsection{Recommendations}

\begin{enumerate}
  \item Key sizes $\ge 2048$ bit (3072 or 4096 for long-term security).
  \item Padding scheme OAEP.
  \item Constant-time implementations.
  \item Regular updates to cryptographic libraries.
\end{enumerate}

\section{Performance Analysis}

\begin{center}
\begin{tabular}{lcc}
\toprule
Operation & Complexity & Typical time (2048 bit) \\
\midrule
Key generation & $O(k^3)$ & 1--10 seconds \\
Encryption & $O(k^2)$ & $<$ 1 ms \\
Decryption & $O(k^3)$ & 10--100 ms \\
\bottomrule
\end{tabular}
\end{center}

\subsection{Optimisations}

\begin{itemize}
  \item Chinese Remainder Theorem for faster decryption.
  \item Windowing methods for exponentiation.
  \item Hardware acceleration.
\end{itemize}

\section{Extensions}

\subsection{Multi-Prime RSA}

$n = p_1 p_2 \cdots p_k$: faster decryption via multi-prime CRT.

\subsection{Blinding}

Protection against timing attacks:
\[
  C' = C \cdot r^e \pmod n, \quad
  M' = (C')^d \pmod n, \quad
  M = M' \cdot r^{-1} \pmod n.
\]

%-------------------------------------------------------------------
\part*{Part II: Period Finding as a Resonance Problem}
%-------------------------------------------------------------------

\section{The arithmetic foundation}

Factorisation of $n=pq$ reduces to period search. For random $a$ with
$\gcd(a,n)=1$ the period $r$ satisfies $a^r \equiv 1 \pmod n$. If $r$
is even and $a^{r/2} \not\equiv -1 \pmod n$ then $\gcd(a^{r/2}\pm1,n)$
yields the factors.

\textbf{[K]} The period $r$ always divides $\lambda(n)=\mathrm{lcm}
(p-1,q-1)$. It is an arithmetic quantity, not determined by a choosable
parameter.

\section{The spectral formulation}

\subsection{Why period finding is a resonance problem}

The sequence $f(x)=a^x \bmod n$ is periodic with period $r$. Its
Fourier transform has sharp maxima at $k/r$. Period search is formally
spectral analysis.

\subsection{Primes as generators}

\textbf{[K]} The resonance structure rests on primes -- measured by the
bicoherence test on 2469 Riemann zeros (Doc.~316, August 2026): coupling
at prime powers ($b^2=0.76$--$0.85$), none at composites
($b^2=0.018$--$0.025$), contrast factor 30--50.

\subsection{What the grid parameter $\sigma$ achieves}

\textbf{[X]} All $\sigma$ values have composite denominators and form
mixed points without their own resonance line. Their function is that of
a sampling grid.

\section{Topology: winding number instead of frequency ratio}

\textbf{[K]} For closed trajectories the winding number $W=\log_2(r)$
decides, not $r$. Pure fifth: $W=0.5850$ (irrational) $\to$ open.
Tempered: $W=7/12$ (rational) $\to$ closed after 12 steps. Closure is
an artefact of rationalisation.

\section{Structural scaling limit}

\subsection{The Weyl obstruction}

\textbf{[K]} The Laplace spectrum of every compact manifold grows as
$N(\lambda)\sim C_d\lambda^{d/2}$; the arithmetic counting density as
$N(T)\sim T\ln T$ -- not a power. A compact resonance space cannot carry
a spectrum of density $T\ln T$.

\textbf{Delimitation.} This obstruction does not apply to Shor's
algorithm. Shor seeks one period of one concrete number (finite task,
$2n+3$ qubits). What limits Shor at large $n$ is the gate count
($\sim n^3$) and error-correction overhead -- a \emph{resource limit}.

\subsection{The precision barrier}

For sharp separation at $r \approx n$ a resolution $\Delta f \approx
f_0/n^2$ is needed -- at RSA sizes below any available number
representation.

\section{Comparison of methods}

\begin{center}
\begin{tabular}{llll}
\toprule
Method & Principle & Closure & Limit \\
\midrule
Trial division & Division & Discrete arithmetic & $O(\sqrt{n})$ \\
Pollard $\rho$ & Cycle detection & Birthday paradox & $O(n^{1/4})$ \\
Resonance filter & Spectral search & Rationalised grid & $O(n)$ \\
Shor's QFT & Quantum interference & Discrete phase &
  resource limit \\
\bottomrule
\end{tabular}
\end{center}

Preparation (modular exponentiation) accounts for approximately
\textbf{95\,\%} of the gate cost; the QFT for only 5\,\%
(Doc.~176, August 2026).

\textbf{[X]} A general exponential quantum advantage is not established.

\section{What the Fourier transform achieves}

\textbf{[K]} The Fourier transform is a \emph{projection}: it translates
a representation without generating new information. Measured on the
Riemann zeros (Doc.~316): peaks at $\ln 2,\ln 3,\ln 5,\ln 7,\ln 11,
\ln 13$ where information is present; no peaks at composite arguments.

\section{Balance}

\begin{center}
\begin{tabular}{ll}
\toprule
Statement & Status \\
\midrule
RSA correctness (Euler, CRT) & \textbf{[K]} proved \\
Period finding is a spectral problem & \textbf{[K]} \\
Primes are the generators & \textbf{[K]} measured by bicoherence \\
Winding number decides closure & \textbf{[K]} \\
Preparation $\approx 95\,\%$ of gate cost & \textbf{[K]} \\
$\sigma$ grids generate resonances & \textbf{[X]} sampling artefact \\
$\xi=4/30000$ contributes to factorisation & \textbf{[X]} geometry \\
Classical period search $O((\log n)^3)$ & \textbf{[X]} refuted \\
Weyl obstruction limits Shor & \textbf{[X]} resource limit \\
Fourier transform generates information & \textbf{[X]} projection \\
\bottomrule
\end{tabular}
\end{center}

\section{Verification}

Scripts in \texttt{2/python/Shor\_Skripte/} (standard library,
assertions, seed 20780458): \texttt{s1\_periodensuche.py},
\texttt{s2\_grenzen.py}, \texttt{s3\_thermik\_zustandsraum.py}.

\section{Relation to the corpus}

Doc.~316 (August 2026): Weyl obstruction, bicoherence test, winding
number, Fourier transform as projection. Revision declared under R75
in Doc.~190.
